\documentclass[12pt,letterpaper]{report}
\usepackage{graphicx}
\usepackage{scrextend}
\usepackage{vmargin}
\usepackage{graphicx}
\usepackage{multirow}
\usepackage[utf8]{inputenc}
\usepackage[spanish]{babel}
\usepackage{multicol}
\usepackage{enumerate}
\usepackage{hyperref}
\usepackage{float}
\usepackage{amsmath, amsthm, amssymb, amsfonts}
\usepackage[usenames]{color}
\definecolor{urlcolor}{rgb}{0,.145,.698}
    \definecolor{linkcolor}{rgb}{.71,0.21,0.01}
    \definecolor{citecolor}{rgb}{.12,.54,.11}
    \definecolor{def}{rgb}{0.00,0.27,0.87}
    \hypersetup{
      breaklinks=true,  % so long urls are correctly broken across lines
      colorlinks=true,
      urlcolor=urlcolor,
      linkcolor=linkcolor,
      citecolor=citecolor,
      }
\parindent=0mm
\pagestyle{empty}
\definecolor{miorange}{rgb}{0.91, 0.43, 0.0}
\begin{document}
\setmargins{2.5cm}      
{1.5cm}                     
{2cm}  
{24cm}                    
{10pt}                          
{1cm}                          
{0pt}                             
{2cm}
Coeficiente de asimetría basado en moda
\begin{equation*}
    A_{p1}= \frac{\mu - MO}{\sigma}
\end{equation*}
Coeficiente de asimetría basado en la mediana
\begin{equation*}
    A_{p2}= \frac{3\left(\mu - Med\right)}{\sigma}
\end{equation*}
Coeficiente de asimetría basado en momentos
\begin{equation*}
    \sqrt{b_1} = \frac{1}{n}\sum_{i=1}^n \left(\frac{x_i-\mu}{\sigma}\right)^3= \frac{1}{n}\sum_{i=1}^n z_i^3 \qquad z_i = \frac{x_i-\mu}{\sigma}
\end{equation*}
Coeficiente de asimetría basado en cumulantes de Fisher
\begin{equation*}
    \gamma_1 = \frac{n}{(n-1)(n-2)\sigma^3} \sum_{i=1}^n (x_i-\mu)^3
\end{equation*}
Coeficiente de asimetría intercuatílico de Bowley y Yule
\begin{equation*}
A_{BY}= \frac{C_1+C_3-2C_2}{C_3-C_1}
\end{equation*}
donde $C_i$ es el i-esimo cuartil.\\
Coeficiente de asimetría percentílico de Kelley
\begin{equation*}
    A_K=\frac{P_{10}+P_{90}-2P_{50}}{P_{90}-P_{10}}
\end{equation*}
donde $P_i$ es el i-esimo percentil.\\
Coeficiente de curtosis basado en momentos de Pearson
\begin{equation*}
    \beta_2 = \frac{1}{n} \sum_{i=1}^n \left(\frac{x_i-\mu}{\sigma}\right)^4
\end{equation*}
Exceso de curtosis de Pearson
\begin{equation*}
    B_2 =n \frac{\sum\limits_{i=1}^n (x_i-\mu)^4}{\sum\limits_{i=1}^n(x_i-\mu)^2} - 3 \left(\frac{n-1}{n+1}\right)
\end{equation*}
Coeficiente de curtosis basado en cumulantes de Fisher
\begin{equation*}
     \gamma_2 = \frac{n(n+1)}{(n-1)(n-2)(n-2)}\sum_{i=1}^n \left(\frac{x_i\mu}{\sigma}\right)^4
\end{equation*}
Exceso de curtosis de Fisher
\begin{equation*}
    \Gamma_2=\frac{n(n+1)}{(n-1)(n-2)(n-3)} \sum_{i=1}^n \left(\frac{x_i-\mu}{\sigma}\right)^4 - 3\frac{(n-1)^2}{(n-2)(n-3)}
\end{equation*}
Coeficiente de curtosis percentílico de Kelley
\begin{equation*}
    CP= \frac{P_{75}-P_{25}}{2(P_{90}-P_{10})}
\end{equation*}
Coeficiente de curtosis percentílico centrado en 0
\begin{equation*}
    CPC= \frac{P_{75}-P_{25}}{2(P_{90}-P_{10})} - 0.263
\end{equation*}
\end{document}